% Chapter 4, Topic _Linear Algebra_ Jim Hefferon
%  http://joshua.smcvt.edu/linearalgebra
%  2020-Apr-08
\topic{Grafika Komputer}
\index{grafika komputer|(}
Topik sebelumnya tentang Geometri Proyektif memberikan model berikut tentang
cara mata kita, atau sebuah kamera, melihat dunia.
\begin{center} 
  \vcenteredhbox{\includegraphics{det/mp/ch4.14}}
\end{center}
Semua titik pada suatu garis yang melalui titik asal diproyeksikan ke tempat
yang sama.

Dalam topik tersebut, kita mendefinisikan bahwa bagi sebarang vektor tak nol
$\vec{v}\in\Re^3$, \definend{titik $p$ pada bidang proyektif}
\index{titik!pada bidang proyektif} yang terkait adalah himpunan
$\set{k\vec{v}\suchthat \text{$k\in\Re$ dan $k\neq 0$}}$.
Himpunan ini terdiri atas vektor-vektor tak nol yang terletak pada garis yang
sama melalui titik asal dengan $\vec{v}$.

Untuk mendeskripsikan sebuah titik proyektif, kita dapat memberikan sebarang
vektor wakil pada garis tersebut.
Karena itu, setiap vektor berikut merepresentasikan titik proyektif yang sama.
\begin{equation*}
  \colvec[r]{1 \\ 2 \\ 3}
  \qquad
  \colvec{1/3 \\ 2/3 \\ 1}
  \qquad
  \colvec[r]{-2 \\ -4 \\ -6}
\end{equation*} 
Masing-masing merupakan \definend{vektor koordinat homogen}
\index{koordinat!homogen}%
\index{vektor koordinat homogen}\index{vektor!koordinat homogen}
bagi titik~$p$.
Dua vektor koordinat homogen (yang menurut definisi tidak nol)
\begin{equation*}
  \tilde{p}_1=\colvec{a_1 \\ b_1 \\ c_1}
  \qquad
  \tilde{p}_2=\colvec{a_2 \\ b_2 \\ c_2}
\end{equation*}
merepresentasikan titik proyektif yang sama jika terdapat faktor skala
$s\neq 0$ sedemikian sehingga $s\tilde{p}_1=\tilde{p}_2$.

Dari tak hingga banyak wakil suatu titik dengan $c\neq0$, kita sering
menggunakan wakil yang komponen ketiganya adalah~$1$.
Hal ini sama dengan memproyeksikan ke bidang $z=1$.
Titik dengan $c=0$ terletak pada garis di tak hingga dan tidak memiliki wakil
pada bagan hingga ini.
\begin{center} 
  \vcenteredhbox{\includegraphics{det/mp/ch4.18}}
\end{center}

Dalam topik ini, kita akan menunjukkan cara menggunakan gagasan-gagasan tersebut
untuk menghasilkan beberapa efek grafika komputer.
Untuk itu, kita ambil gambar sebelumnya dan menggambarnya kembali tanpa bola,
dengan proyektor film di titik asal dan bidang $z=1$ tampak seperti layar
bioskop.
\begin{center} 
  \vcenteredhbox{\includegraphics{det/mp/ch4.60}}
\end{center}
Bagi $c\neq0$, hal ini mengaitkan vektor-vektor dalam ruang tiga dimensi pada
garis abu-abu dengan $p$~pada bidang layar.
\begin{equation*}
  \tilde{p}=\colvec{a \\ b \\ c} 
  \mapsto 
  p=\colvec{x \\ y}=\colvec{a/c \\ b/c}
\end{equation*}

Kita dapat menyesuaikan hal-hal yang telah kita pelajari tentang matriks untuk
melakukan transformasi.
Rotasi\index{rotasi} merupakan salah satu contohnya.
Matriks berikut merotasi bidang~$z=1$ terhadap titik asal sejauh sudut~$\theta$.
\begin{equation*}
  \begin{mat}
    \cos\theta  &-\sin\theta  &0  \\
    \sin\theta  &\cos\theta   &0  \\
    0           &0            &1  
  \end{mat}
  \colvec{x \\ y \\ 1}
  =
  \colvec{\cos\theta\cdot x -\sin\theta\cdot y  \\
         \sin\theta\cdot x +\cos\theta\cdot y  \\ 
         1}
\end{equation*}
Perhatikan bahwa matriks tersebut bekerja pada sebarang vektor koordinat
homogen. Jika kita menerapkan matriks
\begin{equation*}
  \begin{mat}
    \cos\theta  &-\sin\theta  &0  \\
    \sin\theta  &\cos\theta   &0  \\
    0           &0            &1  
  \end{mat}
  \colvec{a \\ b \\ c}
  =
  \colvec{\cos\theta\cdot a -\sin\theta\cdot b  \\
         \sin\theta\cdot a +\cos\theta\cdot b  \\ 
         c}
\end{equation*}
dan kemudian berpindah ke bidang~$z=1$,
\begin{equation*}
  \colvec{\cos\theta\cdot a -\sin\theta\cdot b  \\
         \sin\theta\cdot a +\cos\theta\cdot b  \\ 
         c}
  \mapsto
  \colvec{(\cos\theta\cdot a -\sin\theta\cdot b)/c  \\
         (\sin\theta\cdot a +\cos\theta\cdot b)/c  \\ 
         1}
\end{equation*}
kita memperoleh hasil yang sama seperti jika kita terlebih dahulu berpindah ke
bidang tersebut lalu menerapkan matriks.
\begin{equation*}
  \begin{mat}
    \cos\theta  &-\sin\theta  &0  \\
    \sin\theta  &\cos\theta   &0  \\
    0           &0            &1  
  \end{mat}
  \colvec{a/c \\ b/c \\ 1}
  =
  \colvec{\cos\theta\cdot a/c -\sin\theta\cdot b/c  \\
         \sin\theta\cdot a/c +\cos\theta\cdot b/c  \\ 
         1}
\end{equation*}

Jadi, bekerja dengan koordinat homogen tidak menimbulkan masalah.
Namun, apa keuntungannya?

Operasi translasi dalam grafika komputer,\index{translasi} yaitu menggeser
objek dari satu tempat ke tempat lain, bukan merupakan transformasi linear
karena tidak mempertahankan titik asal.
Namun, jika kita bekerja dengan koordinat homogen, kita dapat menggunakan
matriks.
Matriks berikut mentranslasikan titik-titik pada bidang yang ditinjau sejauh
$t_x$ dalam arah~$x$ dan $t_y$ dalam arah~$y$.
\begin{equation*}
  \begin{mat}
    1  &0  &t_x  \\
    0  &1  &t_y  \\
    0  &0  &1  
  \end{mat}
  \colvec{a \\ b \\ c}
  =
  \colvec{a+t_x\cdot c  \\  
          b+t_y\cdot c  \\ 
          c}
  \mapsto
  \colvec{a/c+t_x  \\  
          b/c+t_y  \\ 
          1}
\end{equation*}
Artinya, pada bidang yang ditinjau, matriks ini menggeser
$\binom{a/c}{b/c}$ ke $\binom{a/c+t_x}{b/c+t_y}$.
Dengan demikian, koordinat homogen memungkinkan kita menggunakan matriks.

Lalu, apa keuntungan menggunakan matriks-matriks ini?
Apa yang kita peroleh dari koordinat tambahan tersebut?
Andaikan kita sedang membuat film dengan grafika komputer.
Pada suatu adegan, kamera bergerak menyamping dan berputar sekaligus.
Setiap titik dalam adegan harus ditranslasikan sekaligus dirotasikan.
Alih-alih meminta komputer melakukan dua operasi pada setiap titik, kita dapat
mengalikan kedua matriks terlebih dahulu. Komputer kemudian hanya menerapkan
satu operasi pada setiap titik, yaitu mengalikannya dengan matriks hasil.
Hal ini sangat mempercepat dan menyederhanakan perhitungan.

Berikut beberapa contoh efek yang dapat kita peroleh.
Kita telah membahas rotasi.
Gambar berikut menunjukkan rotasi sebesar setengah radian.
\begin{center} 
  \vcenteredhbox{\includegraphics{det/mp/ch4.61}}
  \quad$\mapsto$\quad
  \vcenteredhbox{\includegraphics{det/mp/ch4.62}}
\end{center}
Berikut translasi dengan $t_x=1.5$ dan~$t_y=0.5$.
\begin{center} 
  \vcenteredhbox{\includegraphics{det/mp/ch4.61}}
  \quad$\mapsto$\quad
  \vcenteredhbox{\includegraphics{det/mp/ch4.67}}
\end{center}

Berikutnya adalah penskalaan.
Matriks ini mengubah skala objek pada bidang sasaran dengan faktor~$s$ dalam
arah~$x$ dan faktor~$t$ dalam arah~$y$.
\begin{equation*}
  \begin{mat}
    s  &0  &0  \\
    0  &t  &0  \\
    0  &0  &1  
  \end{mat}
  \colvec{a/c \\ b/c \\ 1}
  =
  \colvec{s\cdot a/c  \\  
          t\cdot b/c \\ 
          1}
\end{equation*}
Dalam gambar berikut, kita mengubah skala dalam arah~$x$ dengan faktor $s=2.5$
dan dalam arah~$y$ dengan faktor~$t=0.75$.
\begin{center} 
  \vcenteredhbox{\includegraphics{det/mp/ch4.61}}
  \quad$\mapsto$\quad
  \vcenteredhbox{\includegraphics{det/mp/ch4.63}}
\end{center}

Jika kita mengambil $s=t$, seluruh bentuk berubah skala.
Sebagai contoh, jika kita merangkai bingkai-bingkai dengan $s=t=1.10$, dalam
film tersebut objek akan tampak semakin mendekati kita.
\begin{center} 
  \vcenteredhbox{\includegraphics{det/mp/ch4.61}}
  \quad
  \vcenteredhbox{\includegraphics{det/mp/ch4.68}}
  \quad
  \vcenteredhbox{\includegraphics{det/mp/ch4.69}}
  \quad
  \vcenteredhbox{\includegraphics{det/mp/ch4.70}}
\end{center}

Kita dapat mencerminkan objek.
Matriks berikut mencerminkannya terhadap garis $y=x$.
\begin{equation*}
  \begin{mat}
    0  &1  &0  \\
    1  &0  &0  \\
    0  &0  &1  
  \end{mat}
  \colvec{a/c \\ b/c \\ 1}
  =
  \colvec{b/c  \\  
          a/c \\ 
          1}
\end{equation*}
Garis putus-putus pada gambar berikut adalah $y=x$.
\begin{center} 
  \vcenteredhbox{\includegraphics{det/mp/ch4.61}}
  \quad$\mapsto$\quad
  \vcenteredhbox{\includegraphics{det/mp/ch4.64}}
\end{center}

Matriks berikut mencerminkan objek terhadap $y=-x$.
\begin{equation*}
  \begin{mat}
    0   &-1  &0  \\
    -1  &0  &0  \\
    0   &0  &1  
  \end{mat}
  \colvec{a/c \\ b/c \\ 1}
  =
  \colvec{-b/c  \\  
          -a/c \\ 
          1}
\end{equation*}
Garis putus-putus di bawah adalah $y=-x$.
\begin{center} 
  \vcenteredhbox{\includegraphics{det/mp/ch4.61}}
  \quad$\mapsto$\quad
  \vcenteredhbox{\includegraphics{det/mp/ch4.65}}
\end{center}

Transformasi yang lebih rumit juga dimungkinkan.
Berikut suatu \definend{geseran}.\index{geseran}
\begin{equation*}
  \begin{mat}
    1   &1  &0  \\
    0   &1  &0  \\
    0   &0  &1  
  \end{mat}
  \colvec{a/c \\ b/c \\ 1}
  =
  \colvec{a/c+b/c  \\  
          b/c \\ 
          1}
\end{equation*}
Dalam gambar ini, komponen~$y$ titik-titik tidak berubah, sedangkan pada
komponen~$x$ ditambahkan nilai~$y$.
\begin{center} 
  \vcenteredhbox{\includegraphics{det/mp/ch4.61}}
  \quad$\mapsto$\quad
  \vcenteredhbox{\includegraphics{det/mp/ch4.66}}
\end{center}

Keuntungan utama penggunaan matriks untuk semua operasi ini adalah kita dapat
melakukan hal yang rumit dengan menggabungkan hal-hal sederhana.
Untuk mencerminkan terhadap garis $y=-x+2$, kita dapat menggunakan tiga matriks
untuk mentranslasikan bidang sehingga garis cermin itu berimpit dengan
$y=-x$, mencerminkannya terhadap garis tersebut, lalu mentranslasikannya
kembali.
\begin{equation*}
  \begin{mat}
    1   &0  &0  \\
    0   &1  &2  \\
    0   &0  &1  
  \end{mat}
  \begin{mat}
    0   &-1  &0  \\
    -1  &0   &0  \\
    0   &0   &1  
  \end{mat}
  \begin{mat}
    1   &0  &0  \\
    0   &1  &-2  \\
    0   &0  &1  
  \end{mat}
  \tag{*}
\end{equation*}
(Seperti biasa, tindakan yang dilakukan pertama kali dideskripsikan oleh
matriks paling kanan.
Artinya, matriks paling kanan menggeser semua titik pada bidang yang ditinjau
sejauh $-2$ dalam arah~$y$, matriks tengah mencerminkan terhadap $y=-x$, dan
matriks paling kiri menggeser semua titik kembali.)

Matriks bahkan memungkinkan efek yang lebih rumit.
Berikut matriks bagi \definend{transformasi afin umum},
\index{transformasi afin} dan
\definend{transformasi proyektif umum}.\index{transformasi proyektif}
\begin{equation*}
  \begin{mat}
    d   &e  &f  \\
    g   &h  &i  \\
    0   &0  &1  
  \end{mat}
  \qquad
  \begin{mat}
    d   &e  &f  \\
    g   &h  &i  \\
    j   &k  &\ell
  \end{mat}
\end{equation*}
Untuk transformasi proyektif, matriksnya harus nonsingular.
Namun, deskripsi efek geometris transformasi-transformasi tersebut berada di
luar cakupan kita.

Terdapat pustaka yang sangat luas mengenai grafika komputer, dan aljabar linear
memegang peranan penting di dalamnya.
Salah satu sumber yang sangat baik adalah \cite{Hughes}.
Bidang ini merupakan perpaduan yang indah antara matematika dan seni; lihat
\cite{Disney}.



% ========================================
\begin{exercises}
  \item 
    Hitung hasil kali dalam~($*$).
    \begin{answer}
      Perhitungan langsung memberikan hasil kali berikut.
      \begin{equation*}
        \begin{mat}
          1   &0  &0  \\
          0   &1  &2  \\
          0   &0  &1  
        \end{mat}
        \begin{mat}
          0   &-1  &0  \\
          -1  &0   &0  \\
          0   &0   &1  
        \end{mat}
        \begin{mat}
          1   &0  &0  \\
          0   &1  &-2  \\
          0   &0  &1  
        \end{mat}
        =
        \begin{mat}
         0 &-1 &2 \\
        -1 &0  &2 \\
         0 &0  &1
        \end{mat}
      \end{equation*}
    \end{answer}
% sage: M0 = matrix(QQ, [[1,0,0], [0,1,2], [0,0,1]])
% sage: M0
% [1 0 0]
% [0 1 2]
% [0 0 1]
% sage: M1 = matrix(QQ, [[0,-1,0], [-1,0,0], [0,0,1]])
% sage: M2 = matrix(QQ, [[0,-1,0], [-1,0,0], [0,0,1]])
% sage: M2 = matrix(QQ, [[1,0,0], [0,1,-2], [0,0,1]])
% sage: M0*M1*M2
% [ 0 -1  2]
% [-1  0  2]
% [ 0  0  1]
  \item Cari matriks yang mencerminkan terhadap garis $y=2x$.
    \begin{answer}
      Dengan bekerja dalam $\R^2$, misalkan matriks tersebut adalah $M$.
      Kita memperoleh kedua pemetaan berikut.
      \begin{equation*}
        \colvec{1 \\ 2}\mapsto \colvec{1 \\ 2}
        \qquad
        \colvec{-2 \\ 1}\mapsto \colvec{2 \\ -1}
      \end{equation*}
      (Bagi pemetaan kedua, vektor awal berada pada garis melalui titik asal
      yang tegak lurus terhadap $y=2x$.)
      Kita memiliki
      \begin{equation*}
        M
        \begin{mat}
          1  &-2  \\
          2  &1
        \end{mat}
        =
        \begin{mat}
          1  &2  \\
          2  &-1
        \end{mat}
      \end{equation*}
      Penyelesaian persamaan tersebut memberikan
      \begin{equation*}
        M =
            \begin{mat}
              1  &2  \\
              2  &-1
            \end{mat}
            \begin{mat}
              1  &-2  \\
              2  &1     
            \end{mat}^{-1}
          =
            \begin{mat}
              -3/5  &4/5 \\
              4/5   &3/5
            \end{mat}
      \end{equation*}
      sehingga bagi koordinat homogen, matriksnya adalah
      \begin{equation*}
      \begin{mat}
        -3/5  &4/5  &0  \\
         4/5  &3/5  &0  \\
           0  &0    &1
      \end{mat}
      \end{equation*}
    \end{answer}
  \item Cari matriks yang mencerminkan terhadap garis $y=2x-4$.
    \begin{answer}
      Geser semua titik ke atas sejauh $4$, cerminkan terhadap garis $y=2x$
      menggunakan latihan sebelumnya, lalu geser kembali ke bawah sejauh $4$.
      \begin{equation*}
      \begin{mat}
        1  &0  &0  \\
        0  &1  &-4 \\
        0  &0  &1
      \end{mat}
      \begin{mat}
        -3/5  &4/5  &0  \\
         4/5  &3/5  &0  \\
           0  &0    &1
      \end{mat}
      \begin{mat}
        1  &0  &0  \\
        0  &1  &4 \\
        0  &0  &1
      \end{mat}
      =
      \begin{mat}
        -3/5  &4/5  &16/5  \\
         4/5  &3/5  &-8/5 \\
           0  &0    &1
      \end{mat}
      \end{equation*}
    \end{answer}
% sage: M0 = matrix(QQ, [[1,0,0], [0,1,-4], [0,0,1]])
% sage: M1 = matrix(QQ, [[-3/5,4/5,0], [4/5,3/5,0], [0,0,1]])
% sage: M2 = matrix(QQ, [[1,0,0], [0,1,4], [0,0,1]])
% sage: M0*M1*M2
% [-3/5  4/5 16/5]
% [ 4/5  3/5 -8/5]
% [   0    0    1]
  \item 
    Rotasi dan translasi merupakan operasi kaku.
    Apa matriks untuk rotasi yang diikuti oleh translasi?
    \begin{answer}
    \begin{equation*}
      \begin{mat}
        1  &0  &t_x  \\
        0  &1  &t_y  \\
        0  &0  &1
      \end{mat}
      \begin{mat}
        \cos\theta &-\sin\theta  &0 \\
        \sin\theta &\cos\theta   &0 \\
        0          &0            &1
      \end{mat}
      =
      \begin{mat}
        \cos\theta  &-\sin\theta  &t_x  \\
        \sin\theta  &\cos\theta   &t_y  \\
        0           &0            &1
      \end{mat}
    \end{equation*}
    \end{answer}
  \item Koordinat homogen diperluas ke manipulasi ruang tiga dimensi dengan cara
    yang wajar:~setiap titik proyektif merupakan himpunan vektor tak nol
    berkomponen empat yang satu sama lain merupakan kelipatan skalar.
    Berikan matriks untuk melakukan rotasi terhadap sumbu~$z$ dan matriks untuk
    rotasi terhadap sumbu~$y$.
    \begin{answer}
    Berikut kedua matriks $\nbyn{4}$ tersebut. Matriks kedua menggunakan
    konvensi bahwa sudut positif memutar sumbu~$x$ menuju sumbu~$z$.
    \begin{equation*}
      \begin{mat}
        \cos\theta  &-\sin\theta  &0  &0  \\
        \sin\theta  &\cos\theta   &0  &0  \\
        0           &0            &1  &0  \\
        0           &0            &0  &1  
      \end{mat}
      \qquad
      \begin{mat}
        \cos\theta  &0  &-\sin\theta  &0   \\
        0           &1  &0            &0  \\
        \sin\theta  &0  &\cos\theta   &0   \\
        0           &0  &0            &1  
      \end{mat}
    \end{equation*}    
    \end{answer}
\end{exercises}
\index{grafika komputer|)}
\endinput
